\documentclass[aspectratio=169,11pt]{beamer}
\usepackage[utf8]{inputenc}
\usepackage[T1]{fontenc}
\usepackage[spanish,es-nodecimaldot]{babel}
\usepackage{amsmath,amssymb,booktabs,pgfplots,listings}
\pgfplotsset{compat=1.18}
\usetheme{default}
\usecolortheme{seahorse}
\setbeamertemplate{navigation symbols}{}
\setbeamertemplate{footline}[frame number]
\setbeamerfont{frametitle}{size=\large}
\newcommand{\E}{\mathbb E}
\newcommand{\Var}{\operatorname{Var}}
\newcommand{\Cov}{\operatorname{Cov}}
\newcommand{\K}{\mathcal K}
\newcommand{\idea}[1]{\par\medskip\textbf{Interpretación:} #1}
\lstset{basicstyle=\ttfamily\scriptsize,columns=fullflexible,keepspaces=true,
 escapeinside={(*@}{@*)},
 literate={ℝ}{{\(\mathbb R\)}}1 {≠}{{\(\ne\)}}1 {μ}{{\(\mu\)}}1
 {≤}{{\(\leq\)}}1 {∀}{{\(\forall\)}}1 {∧}{{\(\land\)}}1}
\title{Transparencia, sentimiento y fundamentos}
\subtitle{Una señal pública con ruido y su efecto sobre el precio esperado}
\author{David Julio Olano Silva Nevado}
\date{4 de noviembre de 2026}
\begin{document}
\begin{frame}[plain]
\titlepage
\begin{center}\small\url{https://github.com/davidolano03/ai-project}\end{center}
\end{frame}

\begin{frame}{La pregunta y el resultado}
\large
\textquestiondown{}Una señal pública más precisa hace que el precio responda
más a los fundamentos de una empresa?
\normalsize\medskip
\[
 \K=\frac{\partial\,\E[P\mid\theta,\rho_1,\rho_2]}{\partial\theta}.
\]
Si \(\K=0{,}6\), S/ 1 adicional de noticia produce una respuesta esperada
del precio de S/ 0,60.
\idea{Obtengo una respuesta creciente con precisión pública, manteniendo
fijo el sentimiento y sin aprendizaje adicional desde precios.}
\end{frame}

\begin{frame}{La aplicación: una acción de una empresa cotizada}
\begin{itemize}
\item La empresa es el emisor de la acción.
\item Las instituciones observan la innovación fundamental \(\theta\).
\item Los particulares reciben un reporte público imperfecto.
\end{itemize}
\medskip
Ejemplo ilustrativo: una noticia sobre la demanda futura de sus productos.
\idea{Estudio cómo negocian inversionistas con información desigual.
La divulgación de la empresa es exógena.}
\small\medskip\par
Fuente: Zhang y Zhang (2023), secciones 3 y 6.2.
La observación institucional exacta es una idealización.
\end{frame}

\begin{frame}{Quién conoce la información importa}
\small
\begin{tabular}{p{4.1cm}p{7.4cm}}
\toprule
Situación & Tratamiento en este trabajo\\
\midrule
Solo el emisor conoce el dato
& Los inversionistas lo desconocen. No es la \(\theta\) observada por instituciones.\\[0.15cm]
Algunos inversionistas conocen mejor los fundamentos
& Es la asimetría estudiada, idealizada como observación de \(\theta\).\\[0.15cm]
Un directivo transmite un dato a un agente
& El modelo no especifica esa transmisión ni los incentivos del emisor.\\
\bottomrule
\end{tabular}
\normalsize
\idea{Información privada describe quién sabe qué. Su origen necesita
una hipótesis adicional.}
\end{frame}

\begin{frame}{El pago futuro y la secuencia}
\[
 \underbrace{F}_{\text{pago final}}
 =\underbrace{P_0}_{\text{referencia conocida}}
 +\underbrace{\theta}_{\text{fundamento}}
 +\underbrace{\epsilon}_{\text{sorpresa final}}.
\]
\(\theta\sim N(0,s)\), \(\epsilon\sim N(0,e)\), independientes,
\(s>0,\ e>0\).
\medskip\par
Primero se observan señales y se negocia. Después se realiza \(\epsilon\).
\idea{\(\theta=10\) añade S/ 10 al pago esperado.
La sorpresa final todavía puede cambiar el pago realizado.}
\small\medskip\par
\(F,P_0,\theta,\epsilon\): dinero por acción.
\(s,e\): unidades monetarias al cuadrado por acción al cuadrado.
\end{frame}

\begin{frame}{El punto de partida que voy a corregir}
En el modelo básico del paper, las instituciones observan \(\theta\).
Los particulares no reciben información pública sobre su realización.
\[
 \text{Particular racional:}\qquad
 \E[F]=P_0,\quad \Var(F)=e+s.
\]
La extensión publicada introduce
\[
 \kappa=\mu\theta,\qquad v_3^I=e+\mu^2s.
\]
\idea{La sección 6.2 intenta representar divulgación parcial.
Revisaré si esas ecuaciones conservan información oculta.}
\small\medskip\par
Fuente: ecuaciones (15), (39)--(40), pp. 5 y 15.
\end{frame}

\begin{frame}{Escalar un dato exacto permite reconstruirlo}
\[
 \kappa=\mu\theta,\quad \mu>0
 \quad\Longrightarrow\quad \theta=\frac{\kappa}{\mu}.
\]
Si el público conoce \(\mu=0{,}1\) y observa \(\kappa=1\),
recupera \(\theta=10\).
\idea{La señal necesita ruido o un componente desconocido para dejar
incertidumbre real. Reducir su tamaño numérico no basta.}
\small\medskip\par
Auditoría de la ecuación (39), con parámetro conocido.
Lean comprueba la identidad algebraica.
\end{frame}

\begin{frame}{La varianza publicada no conserva los extremos}
La varianza mide cuánto pueden separarse los resultados del pronóstico.
\[
 v_3^I=e+\mu^2s.
\]
\begin{center}
\begin{tabular}{lll}
\toprule
Información & Fórmula publicada & Riesgo restante coherente\\
\midrule
Ninguna, \(\mu=0\) & \(e\) & \(e+s\)\\
Completa, \(\mu=1\) & \(e+s\) & \(e\)\\
\bottomrule
\end{tabular}
\end{center}
\idea{Una noticia ya observada deja de ser una incógnita.
La sorpresa final sigue siendo incierta.}
\small\medskip\par
La columna derecha fija los extremos. La fórmula intermedia requiere
construir una estructura de información.
\end{frame}

\begin{frame}{El contraejemplo compara la extensión con el básico}
\begin{center}
\begin{tabular}{lc}
\toprule
Escenario & Respuesta por S/ 1 de fundamento\\
\midrule
Básico, sin información pública & S/ \(0{,}643\)\\
Extensión impresa, \(\mu=0\) & S/ \(0{,}500\)\\
Extensión impresa, \(\mu=0{,}1\) & S/ \(0{,}552\)\\
\bottomrule
\end{tabular}
\end{center}
La proposición 5(ii) afirma mejora frente al básico, pero \(0{,}552<0{,}643\).
\idea{Un caso admisible cuestiona la afirmación universal impresa.
La comparación relevante es la primera fila con la tercera.}
\small\medskip\par
\(\lambda_i=1/4,\ \gamma=R=1\), sesgos y estados cero.
\(e=s=\tau_1^2=\tau_2^2=1\). Lean verifica las fracciones exactas.
\end{frame}

\begin{frame}{Nuestra señal pública contiene ruido}
\[
 \underbrace{y}_{\text{reporte observado}}
 =\underbrace{\theta}_{\text{fundamento}}
 +\underbrace{\eta}_{\text{ruido del reporte}},
 \qquad \eta\sim N(0,t).
\]
\(\eta\) es independiente de \(\theta,\epsilon\), con \(t>0\).
\medskip\par
\(t\) mide imprecisión del reporte: cuanto menor es, más útil resulta.
\idea{Observar el reporte no permite separar exactamente fundamento y ruido.}
\small\medskip\par
\(y,\eta\) se miden en dinero por acción. La institución observa además
\(\theta\). El reporte resume información económicamente relevante.
\end{frame}

\begin{frame}{Bayes: combinar el conocimiento previo con el reporte}
Para un posible fundamento \(x\) y un dato observado \(y\):
\[
 f(x\mid y)\propto f(y\mid x)f(x)
 \propto\exp\left\{-\frac12
 \left[\frac{x^2}{s}+\frac{(y-x)^2}{t}\right]\right\}.
\]
\[
 \frac{x^2}{s}+\frac{(y-x)^2}{t}
 =\left(\frac1s+\frac1t\right)x^2-\frac{2y}{t}x+\frac{y^2}{t}.
\]
\idea{La distribución previa aporta \(x^2/s\).
La distancia entre el fundamento posible y el reporte aporta \((y-x)^2/t\).}
\end{frame}

\begin{frame}{Completar el cuadrado entrega el pronóstico}
\[
 \frac{x^2}{s}+\frac{(y-x)^2}{t}
 =\frac{s+t}{st}\left(x-\frac{s}{s+t}y\right)^2+\frac{y^2}{s+t}.
\]
El último término es constante respecto de \(x\). Reconozco una normal:
\[
 \theta\mid y\sim
 N\left(\underbrace{\frac{s}{s+t}y}_{\text{media posterior}},
 \underbrace{\frac{st}{s+t}}_{\text{varianza posterior}}\right).
\]
\idea{El reporte cambia el pronóstico y reduce la incertidumbre previa \(s\).}
\end{frame}

\begin{frame}{Definir \(\mu\), \(\kappa\) y \(\xi\) después de Bayes}
\[
 \underbrace{\mu=\frac{s}{s+t}}_{\text{peso del reporte}},\qquad
 \underbrace{\kappa=\mu y}_{\text{pronóstico público}},\qquad
 \underbrace{\xi=\theta-\kappa}_{\text{error todavía oculto}}.
\]
\(\theta=\kappa+\xi\). Con menor ruido \(t\), el público asigna más peso
al reporte y \(\mu\) aumenta.
\idea{\(\kappa\) resume lo aprendido sobre el fundamento.
Dividir por \(\mu\) solo recupera el reporte ruidoso \(y\).}
\small\medskip\par
Esta construcción justifica las proporciones de varianza de la propuesta.
\end{frame}

\begin{frame}{Derivación de la varianza pública}
Como \(\theta\) y \(\eta\) son independientes,
\[
 \Var(y)=\Var(\theta+\eta)=s+t.
\]
Escalar una variable multiplica su varianza por el cuadrado del factor:
\[
 \Var(\kappa)=\Var(\mu y)=\mu^2(s+t)
 =\frac{s^2}{(s+t)^2}(s+t)
 =\frac{s^2}{s+t}=\boxed{\mu s}.
\]
\idea{La proporción surge del pronóstico de Bayes y de la varianza del
reporte. No se obtiene multiplicando directamente \(\theta\) por \(\mu\).}
\end{frame}

\begin{frame}{Derivación de la varianza oculta}
Primero reescribo el residuo:
\[
 \xi=\theta-\mu(\theta+\eta)=(1-\mu)\theta-\mu\eta.
\]
La independencia elimina el término cruzado:
\[
 \Var(\xi)=(1-\mu)^2s+\mu^2t
 =\frac{t^2s+s^2t}{(s+t)^2}
 =\frac{st}{s+t}=\boxed{(1-\mu)s}.
\]
\idea{El residuo combina el fundamento no explicado y el ruido del pronóstico.}
\end{frame}

\begin{frame}{Por qué el componente público no revela el residuo}
\[
 \Cov(\kappa,\xi)
 =\mu(1-\mu)s-\mu^2t
 =\mu[(1-\mu)s-\mu t]=0.
\]
Ambos componentes son conjuntamente normales.
En ese caso, covarianza cero implica independencia.
\[
 \E[\theta\mid\kappa]=\kappa,\qquad
 \Var(\theta\mid\kappa)=(1-\mu)s.
\]
\idea{Después de observar \(\kappa\), sigue faltando \(\xi\).
Su incertidumbre disminuye al aumentar la precisión pública.}
\end{frame}

\begin{frame}{Un ejemplo que distingue promedio e incertidumbre}
Si \(s=t=100\), entonces \(\mu=1/2\). El reporte es \(y=8\):
\[
 \kappa=\tfrac12(8)=4,\qquad
 \Var(\xi)=\tfrac12(100)=50.
\]
Con \(P_0=100\) y \(e=4\):
\[
 \E[F\mid y]=104,\qquad \Var(F\mid y)=50+4=54.
\]
\idea{Espero 104, pero puedo equivocarme.
La parte oculta y la sorpresa final explican el riesgo que permanece.}
\small\medskip\par
Varianza 50 implica desviación estándar \(\sqrt{50}\simeq7{,}071\).
Las varianzas se expresan en unidades monetarias al cuadrado.
\end{frame}

\begin{frame}{Los extremos ahora representan lo que queremos}
\begin{center}
\begin{tabular}{lccc}
\toprule
Precisión pública & \(\mu=0\) & \(\mu=0{,}5\) & \(\mu=1\)\\
\midrule
Varianza pública, \(s=100\) & 0 & 50 & 100\\
Varianza oculta & 100 & 50 & 0\\
\bottomrule
\end{tabular}
\end{center}
En cero defino ausencia de señal: \(\kappa=0,\xi=\theta\).
En uno defino señal exacta: \(\kappa=\theta,\xi=0\).
\idea{La variación total sigue siendo 100.
Una parte cada vez mayor pasa a ser observable.}
\small\medskip\par
El 50\% se refiere a varianza, no a la mitad del valor realizado de \(\theta\).
\end{frame}

\begin{frame}{Las creencias de los cuatro grupos}
\small
\begin{center}
\begin{tabular}{llll}
\toprule
Grupo & Media \(m_i\) & Varianza \(v_i\) & Aversión \(g_i\)\\
\midrule
Institución racional & \(P_0+\theta\) & \(e\) & \(\gamma\)\\[0.1cm]
Institución sentimental & \(P_0+\theta+\delta_1\) & \(e+\tau_1^2\) & \(\gamma e^{-b\rho_1}\)\\[0.1cm]
Particular racional & \(P_0+\kappa\) & \(e+(1-\mu)s\) & \(\gamma\)\\[0.1cm]
Particular sentimental & \(P_0+\kappa+\delta_2\) & \(e+(1-\mu)s+\tau_2^2\) & \(\gamma e^{-b\rho_2}\)\\
\bottomrule
\end{tabular}
\end{center}
\normalsize
\(\delta_j\): sesgo en el pronóstico. \(\tau_j^2\): varianza subjetiva adicional.
\(\rho_j\): estado sentimental. \(b>0\): sensibilidad de aversión al estado.
\idea{Todos reciben el mismo pago real. La tabla describe sus creencias.}
\end{frame}

\begin{frame}{Cómo se obtiene cada varianza del pago}
Para el particular racional, \(\kappa\) ya es conocido:
\[
 F=P_0+\kappa+\xi+\epsilon.
\]
\[
 v_3=\Var(\xi+\epsilon\mid\kappa)
 =(1-\mu)s+e+
 \underbrace{2\Cov(\xi,\epsilon\mid\kappa)}_{0}.
\]
Para una institución, \(\theta\) ya es conocido y \(v_1=e\).
Un error sentimental independiente añade \(\tau_j^2\).
\idea{Las independencias justifican sumar varianzas.
Un sesgo constante cambia la media, sin añadir varianza.}
\end{frame}

\begin{frame}{CARA: cómo el agente valora la riqueza}
\[
 U_i(W)=-e^{-g_iW},\qquad g_i>0.
\]
\[
 U_i'(W)=g_i e^{-g_iW}>0,\qquad
 U_i''(W)=-g_i^2e^{-g_iW}<0.
\]
\[
 -\frac{U_i''(W)}{U_i'(W)}=g_i.
\]
\idea{El agente prefiere más riqueza y evita riesgo.
Su aversión absoluta es constante respecto de su riqueza.}
\small\medskip\par
\(g_i\) tiene unidades de inverso de dinero y permanece fijo al elegir.
\end{frame}

\begin{frame}{Comprar acciones modifica media y varianza de riqueza}
\[
 W_i=(w_{0i}-P\omega_i)R+\omega_i F
 =\underbrace{Rw_{0i}}_{C_{0i}}+\omega_i(F-RP).
\]
\[
 M_i=\E_i[W_i]=C_{0i}+\omega_i(m_i-RP),
\]
\[
 V_i=\Var_i(W_i)
 =\omega_i^2\Var_i(F-RP)=\boxed{\omega_i^2v_i}.
\]
\idea{Duplicar la cantidad de acciones cuadruplica la varianza de riqueza.
El precio pagado es conocido durante la elección.}
\small\medskip\par
\(\omega_i\): cantidad de acciones. \(R=1+r_f\): factor seguro.
\end{frame}

\begin{frame}{La normalidad conecta CARA con media-varianza}
Para riqueza normal, \(\E[e^{aW}]=e^{aM+a^2V/2}\).
\[
 \E_i[U_i(W_i)]
 =-\exp\left(-g_iM_i+\frac{g_i^2V_i}{2}\right)
 =-\exp\left[-g_i\left(M_i-\frac{g_iV_i}{2}\right)\right].
\]
La transformación \(-e^{-g_i z}\) es creciente. Maximizar utilidad
equivale a maximizar
\[
 CE_i=C_{0i}+\omega_i(m_i-RP)-\frac{g_iv_i}{2}\omega_i^2.
\]
\idea{El equivalente cierto suma beneficio esperado y resta una
penalización por riesgo. La integral completa está en el anexo A1.}
\end{frame}

\begin{frame}{Derivar la CPO, término por término}
\[
 CE_i=C_{0i}+\omega_i(m_i-RP)-\tfrac12g_iv_i\omega_i^2.
\]
\[
 \frac{\partial CE_i}{\partial\omega_i}
 =\underbrace{0}_{\text{constante}}
 +\underbrace{(m_i-RP)}_{\text{término lineal}}
 -\underbrace{g_iv_i\omega_i}_{\text{derivada del cuadrado}}=0.
\]
\[
 \boxed{\omega_i^*=\frac{m_i-RP}{g_iv_i}},\qquad
 \frac{\partial^2 CE_i}{\partial\omega_i^2}=-g_iv_i<0.
\]
\idea{El óptimo equilibra beneficio marginal y costo marginal del riesgo.
La concavidad garantiza un máximo único.}
\end{frame}

\begin{frame}{La demanda del particular racional}
\[
 \omega_3^*=
 \frac{\overbrace{P_0+\kappa}^{\text{pago esperado}}-
       \overbrace{RP}^{\text{costo financiado}}}
 {\underbrace{\gamma}_{\text{aversión}}
  \underbrace{[e+(1-\mu)s]}_{\text{varianza que permanece}}}.
\]
\idea{La información modifica cuánto espera recibir y cuánto riesgo cree
estar tomando. Ambos canales afectan su demanda.}
\medskip\par
La fórmula permite posiciones negativas. El resultado usa esa libertad
de negociar, masas positivas y riesgo final \(e>0\).
\end{frame}

\begin{frame}{El precio iguala demanda agregada y oferta}
\[
 1=\sum_i\lambda_i\omega_i^*
 =\sum_i\frac{\lambda_i}{g_iv_i}(m_i-RP).
\]
Defino \(h_i=\lambda_i/(g_iv_i)\) y \(H=\sum_i h_i\). Entonces
\[
 1=\sum_i h_im_i-RPH,\qquad
 \boxed{P=\frac{\sum_i h_im_i-1}{RH}}.
\]
\idea{Los grupos con mayor masa, menor aversión o menor varianza percibida
tienen mayor peso en la respuesta agregada.}
\small\medskip\par
\(\lambda_i>0\) suman uno. La oferta normalizada es una acción.
\end{frame}

\begin{frame}{De la señal realizada al precio esperado}
Con \(a=h_1+h_2\), \(B=h_3+h_4\) y \(H=a+B\):
\[
 P=\frac{HP_0+a\theta+B\kappa+h_2\delta_1+h_4\delta_2-1}{RH}.
\]
El ruido público tiene media cero:
\[
 \E[\kappa\mid\theta]=\E[\mu(\theta+\eta)\mid\theta]=\mu\theta.
\]
Por tanto,
\[
 \bar P=C(\mu)+\underbrace{\frac{a+\mu B}{R(a+B)}}_{\K(\mu)}\theta.
\]
\idea{\(\K\) mide sensibilidad esperada. El promedio elimina la realización
particular del ruido del reporte.}
\end{frame}

\begin{frame}{La proposición propia y todas sus condiciones}
Si \(0\leq x<y\leq1\), entonces
\[
 \boxed{\K(y)>\K(x)},\qquad 0<\K(\mu)\leq1/R,\quad \K(1)=1/R.
\]
\begin{itemize}
\item Cuatro masas positivas, \(R>0\), \(s>0\), \(e>0\).
\item Estados, sesgos, aversiones y varianzas sentimentales fijos.
\item Ruido público independiente y aprendizaje desde la señal especificada.
\item Oferta fija y posiciones reales sin restricciones.
\end{itemize}
\idea{La conclusión trata de cuánto responde el precio esperado a una noticia.}
\end{frame}

\begin{frame}{Por qué la sensibilidad aumenta}
Sean \(d_j=e_j+(1-\mu)s\), \(B=\sum_{j=3}^4n_j/d_j\),
con \(n_j,e_j>0\). Las instituciones mantienen \(a\) fijo.
\[
 \K'=\frac{B^2+a[B-(1-\mu)B']}{R(a+B)^2}.
\]
Para cada grupo particular:
\[
 \frac{n_j}{d_j}-(1-\mu)\frac{n_js}{d_j^2}
 =\frac{n_je_j}{d_j^2}>0.
\]
\[
 \boxed{\K'=\frac{B^2+a\sum_j n_je_j/d_j^2}{R(a+B)^2}>0}.
\]
\idea{El aumento de respuesta al fundamento domina el cambio de pesos
cuando mantenemos fijo el sentimiento.}
\end{frame}

\begin{frame}{La corrección recupera el básico y aumenta con precisión}
\begin{center}
\resizebox{0.83\textwidth}{!}{% Generated from checked data. No screenshots of published figures.
\begin{tikzpicture}
\begin{axis}[
 width=10.6cm,height=5.7cm,
 xmin=0,xmax=1,ymin=0.45,ymax=1.03,
 xlabel={Transparencia $\mu$},ylabel={Sensibilidad informativa $\mathcal K$},
 tick label style={font=\small},label style={font=\small},
 legend style={font=\scriptsize,draw=none,at={(0.02,0.98)},anchor=north west},
 grid=major,grid style={gray!15},
]
\addplot[blue!70!black,very thick] table[x=mu,y=corrected,col sep=comma]{../code/output/information.csv};
\addlegendentry{Corrección: precio esperado}
\addplot[orange!85!black,thick] table[x=mu,y=printed,col sep=comma]{../code/output/information.csv};
\addlegendentry{Fórmula publicada}
\addplot[black,dashed] table[x=mu,y=baseline,col sep=comma]{../code/output/information.csv};
\addlegendentry{Modelo básico}
\end{axis}
\end{tikzpicture}
}
\end{center}
\small
En \(\mu=0{,}1\): publicado \(0{,}551866\), básico \(0{,}642857\),
corregido \(0{,}669346\). Parámetros del contraejemplo.
\end{frame}

\begin{frame}{Qué comprueba el código}
\begin{itemize}
\item Fracciones exactas del contraejemplo y de la corrección.
\item Posterior, varianzas y covarianza mediante identidades racionales.
\item CPO, máximo y vaciado de mercado en cuatro grupos.
\item 5\,760 comparaciones exactas de sensibilidad y 60 mercados.
\item Momentos gaussianos simulados con semilla fija e integral CARA numérica.
\end{itemize}
\idea{Las simulaciones permiten detectar errores de implementación.
La demostración general depende de los supuestos escritos.}
\small\medskip\par
Un único comando: \texttt{python code/verify.py}.
\end{frame}

\begin{frame}[fragile]{Lean comprueba pasos algebraicos identificables}
Ejemplo: la proporción pública sale de una identidad, no de una hipótesis
que ya afirme el resultado.
\begin{lstlisting}
theorem public_variance (s t : (*@ \(\mathbb R\) @*)) (h : s + t (*@ \(\ne\) @*) 0) :
    (s / (s + t)) ^ 2 * (s + t) =
      (s / (s + t)) * s := by
  field_simp
\end{lstlisting}
\small
Además, \texttt{four\_group\_information\_increases} prueba la
monotonicidad del coeficiente de cuatro grupos.
\normalsize
\idea{Bayes, normalidad conjunta e integral CARA tienen derivación escrita.
Su formalización probabilística completa permanece abierta.}
\end{frame}

\begin{frame}{Dónde corregí la explicación de la IA}
\begin{center}\small
\begin{tabular}{p{4.6cm}p{6.7cm}}
\toprule
Afirmación inicial & Veredicto y comprobación\\
\midrule
Asignar \(\mu s\) y \((1-\mu)s\) basta como derivación
& Era un supuesto. Ahora lo derivo desde \(y=\theta+\eta\) y Bayes.\\[0.15cm]
Media-varianza es la utilidad esperada CARA
& Es un equivalente cierto. La integral y la transformación creciente
justifican la misma elección.\\[0.15cm]
Más transparencia siempre mejora el precio
& El resultado demostrado es sensibilidad esperada bajo condiciones fijas.\\
\bottomrule
\end{tabular}
\end{center}
\normalsize
El paso que debo poder reproducir a mano es la expansión de \(\Var(\xi)\)
y su covarianza con \(\kappa\). La guía está en \texttt{hand/}.
\end{frame}

\begin{frame}{El límite principal: aprender del precio cambiaría el problema}
Si todos conocieran pesos, sesgos y oferta fija, podrían despejar:
\[
 \theta=\frac{RHP-HP_0-B\kappa-h_2\delta_1-h_4\delta_2+1}{a}.
\]
Por eso mantengo explícita la restricción de no inferir información
adicional del precio.
\idea{Resolver un mercado con esa inferencia exige un equilibrio informativo
distinto, por ejemplo con oferta aleatoria o señales institucionales imperfectas.}
\small\medskip\par
Tampoco se demuestra una ganancia extraordinaria institucional:
eso requiere una medida de rendimiento, costos y referencia.
\end{frame}

\begin{frame}{Preguntas para la discusión}
\large
\begin{enumerate}
\item \textquestiondown{}Qué interpretación económica defiende mejor
la ausencia de aprendizaje adicional desde precios?
\item \textquestiondown{}Conviene extender primero el modelo con oferta
aleatoria o con información institucional imperfecta?
\item \textquestiondown{}Qué reacción del sentimiento a la precisión pública
podría cambiar el resultado?
\end{enumerate}
\normalsize\medskip
El paper desarrolla cuatro proposiciones. La prueba Lean actual certifica
su parte algebraica y declara los pasos probabilísticos pendientes.
\end{frame}

\appendix
\begin{frame}{A1. La integral normal detrás de CARA}
Para \(W\sim N(M,V)\), \(V>0\):
\[
 aw-\frac{(w-M)^2}{2V}
 =-\frac{[w-(M+aV)]^2}{2V}+aM+\frac{a^2V}{2}.
\]
\[
 \E[e^{aW}]
 =e^{aM+a^2V/2}
 \underbrace{\int_{\mathbb R}
 \frac{e^{-[w-(M+aV)]^2/(2V)}}{\sqrt{2\pi V}}\,dw}_{1}
 =e^{aM+a^2V/2}.
\]
\idea{Completar el cuadrado deja una densidad que integra uno.
Sustituyo \(a=-g_i\). Si \(V=0\), la riqueza es constante.}
\end{frame}

\begin{frame}{A2. Los términos cruzados de la varianza oculta}
\[
 \Var(A X+B Y)
 =A^2\Var(X)+B^2\Var(Y)+2AB\Cov(X,Y).
\]
Con \(A=1-\mu\), \(B=-\mu\), \(X=\theta\), \(Y=\eta\):
\[
 \Var(\xi)=(1-\mu)^2s+\mu^2t
 -2\mu(1-\mu)\underbrace{\Cov(\theta,\eta)}_0.
\]
\[
 (1-\mu)^2s+\mu^2t
 =\frac{t^2s+s^2t}{(s+t)^2}
 =\frac{st}{s+t}.
\]
\idea{Independencia del ruido y el fundamento elimina la covarianza.}
\end{frame}

\begin{frame}{A3. La CPO da un máximo global}
\[
 CE(\omega^*)-CE(\omega)
 =\frac{gv}{2}(\omega-\omega^*)^2\geq0.
\]
\(g,v>0\), por lo que la igualdad solo ocurre si \(\omega=\omega^*\).
\idea{La posición de la CPO domina cualquier otra posición real.
Lean verifica esta identidad y la desigualdad.}
\small\medskip\par
Si se prohibiera vender en corto, habría que resolver el problema
restringido; la demanda sería el máximo entre cero y la solución interior.
\end{frame}

\begin{frame}{A4. Una diferencia finita evita aproximar derivadas}
Con \(d_j(z)=e_j+(1-z)s\) y \(B_z=\sum_j n_j/d_j(z)\):
\[
 \K(y)-\K(x)=
 \frac{(y-x)\left[B_xB_y+
 a\sum_j\frac{n_je_j}{d_j(x)d_j(y)}\right]}
 {R(a+B_x)(a+B_y)}.
\]
Todos los términos son positivos para \(x<y\leq1\).
\idea{El incremento es exacto y cubre los extremos.
El código comprueba la identidad con fracciones.}
\end{frame}

\begin{frame}{A5. Qué queda para completar el paper del curso}
\begin{itemize}
\item Formalización probabilística de Bayes, gaussianas e integral CARA.
\item Procedimiento integral de AppliedModelingLib sobre un commit fijado.
\item Apéndice manuscrito personal con todas las manipulaciones.
\end{itemize}
\medskip
Los archivos actuales identifican esas tareas pendientes.
\idea{La presentación final es un avance para recibir crítica.
Las instrucciones actuales fijan el paper final para el 26 de noviembre.}
\small\medskip\par
Fuente: issue 7 del curso, secciones 4--6, consultado el 2 de noviembre.
\end{frame}

\begin{frame}{A6. Fuentes y ubicación de resultados}
\small
Zhang y Zhang (2023), \emph{Information asymmetry, sentiment interactions,
and asset price}. NAJEF 67, 101920.
\par\medskip
Modelo: pp. 3--6. Transparencia: sección 6.2, pp. 15--16.
Ecuaciones (39)--(44) y proposición 5(ii).
\par\medskip
Di Maggio y Pagano (2018), \emph{Financial Disclosure and Market
Transparency with Costly Information Processing}. Review of Finance 22(1),
117--153. Publicación en línea: 2017.
\par\medskip
Grossman y Stiglitz (1980), \emph{On the Impossibility of Informationally
Efficient Markets}. AER 70(3), 393--408.
\par\medskip
Derivaciones y condiciones: \texttt{paper/paper.pdf}.
Verificación: \texttt{code/verify.py}. Pruebas: \texttt{lean/preliminary/}.
\end{frame}
\end{document}
